\section{Scope and Conclusion}
\label{sec:discussion}
\label{sec:conclusion}

The evidence concerns the recorded SAGA case and tested policies, not
arbitrary policies, current upstream vulnerability status, or prevalence
across deployments. The finite specification evaluator is not
theorem-proved. Formal results depend on symbolic role mapping, transport,
and handoff assumptions; they do not establish token expiry, quota
consumption, revocation, policy updates, session isolation, or injective
agreement. Consumer validation uses controlled registration and transport
fixtures and stops before message acceptance. Neither the symbolic witness
nor the consumer experiment establishes a deployed-service exploit,
Python/ProVerif refinement, or a verified production repair.

The recorded rule-order inconsistency separates intended policy permission
from the matcher's observed decision. Conditional symbolic analysis admits
message acceptance with prior token issuance, while the independent
consumer experiment observes propagation through real authorization logic
to token issuance and storage. Authentication and credential provenance
therefore do not substitute for decision correctness in this case.
Future work may investigate stronger implementation-to-model correspondence
and additional systems.
